\chapter{Introduction}\label{ch:intro}

miniVivisection extends the debugger of Rhino~8's Script Editor, and it brings that debugger into
Grasshopper. When a script stops at a breakpoint, you see what its variables hold \emph{as geometry in
Rhino's viewports}, not only as lines of text. You can change that geometry with Rhino's own tools and
write it back into the running script. The values appear where you read the code: in the tooltips, and
at the ends of the lines. Rhino stays usable the whole time, so you can select, measure and run commands
while the script waits.

\section{Why it exists}

Scripts that make geometry are hard to debug with a text debugger. A list of four
\code{Point3d} values can look right number by number and still describe a rectangle whose corners
are in the wrong order. On the screen, the mistake is obvious: the outline crosses itself
(the cover picture). Rhino's Script Editor has a good debugger. It shows values as text, though, and
while a script is paused Rhino is locked in a \ui{Select Action} prompt: you can't select an object,
measure a distance or look at the model from another side.

Other problems come up again and again in teaching and practice:
\begin{itemize}
  \item A script started from Rhino's command line hangs in a loop, and the only way out is to kill
        Rhino and lose the model.
  \item A Grasshopper script component fails for three items out of a hundred, and Grasshopper says
        only \emph{“'NoneType' object is not iterable”}, not which item, branch or run.
  \item Notes about a definition sit in Panels as centered plain text, and scripts kept in Panels have
        no line numbers, so an error's line number points nowhere.
\end{itemize}

miniVivisection addresses each of these, inside the tools you already use. It adds nothing you have to
learn before you can start: set a breakpoint, press \key{F5}, and the values are drawn.

\section{What it does, in six ideas}

\begin{description}
  \item[See.] A paused script's values are drawn in Rhino: points, curves, surfaces, meshes, planes,
    vectors and lists of them, each with a name. The code shows its values in tooltips and at the ends
    of the lines, and what changed since the last step (\cref{ch:watch,ch:code}).
  \item[Change and keep.] Turn a value into Rhino objects, edit them with the Gumball or control
    points, and the paused script continues with your edit. Or bake what was computed into the
    document before the script crashes (\cref{ch:edit}).
  \item[Stop where you want.] Breakpoints with conditions, hit counts and log messages. \keys{Ctrl+Alt+B}
    pauses a running script at its next line, even one that holds Rhino's UI thread, or stops a
    plain run and records where it was (\cref{ch:breakpoints,ch:running}).
  \item[Ask the paused script.] A Python prompt runs in the paused frame and can change the script's
    variables. A PowerShell prompt and a chat with Claude work on the same paused script
    (\cref{ch:prompts}).
  \item[Debug Grasshopper.] Script components pause at breakpoints and \keys{Ctrl+Alt+B}. Failing runs
    are named by branch and item, and can be debugged one by one. The canvas stays usable while a
    component is paused (\cref{ch:ghdebug,ch:ghpaused}).
  \item[Document in Markdown.] Panels show scripts as code and notes as rendered Markdown. Markdown
    becomes Rhino text, and links lead from a note to a component, a Rhino object, a named view,
    another file or a script (\cref{ch:panels,ch:markdowntext,ch:links}).
\end{description}

\section{Requirements}

\begin{xltabular}{\linewidth}{@{}P{0.28\linewidth}L@{}}
\toprule
\ui{Rhino} & Rhino~8 for Windows, with its Script Editor (RhinoCode). Rhino~8 runs on .NET; IronPython~2
             scripts need the .NET~8 runtime (Rhino's \code{SetDotNetRuntime}). \\
\ui{Python 3 scripts} & Get every feature. This is what miniVivisection is made for. \\
\ui{IronPython 2 scripts} & (\code{\#! python2}) Get the debugger features, drawing of plain variable names and
             conditional breakpoints; not expressions, inline values, editing in Rhino or the Python
             prompt in the paused frame (\cref{ch:limits}). \\
\ui{C\# scripts} & Watches of plain variable names and conditional breakpoints (a small C\# evaluator);
             largely untested. \\
\ui{Grasshopper} & Rhino~8's Python~3 script components. The legacy GhPython (IronPython) component is not covered. \\
\ui{V°Claude} (optional) & Claude Code, the \code{claude} command line program, installed and signed in.
             RhinoAI (McNeel's Rhino MCP Platform) lets Claude work in the model and in Grasshopper. \\
\ui{V°PowerShell} & PowerShell~7 (\code{pwsh}) or Windows PowerShell. \\
\bottomrule
\end{xltabular}

\begin{caution}[Experimental]
miniVivisection~\mvversion\ is experimental. The Script Editor has no extension interface, so the plug-in
reaches into its internal code. A Rhino service release can change that code and switch parts of the
plug-in off; it then does nothing rather than fail. It was called ScriptEditorWatchViz until October 2026.
\end{caution}

\section{How this manual is organized}

\begin{description}
  \item[Part I, Getting to know miniVivisection:] This introduction, all features at a glance
    (\cref{ch:features}), the concepts the features are built on (\cref{ch:concepts}) and a first session
    (\cref{ch:start}). Read \cref{ch:concepts} once: most questions about why something works the way it
    does are answered there.
  \item[Part II, In the Script Editor:] Every feature for scripts, from drawing values to threads and the
    prompts.
  \item[Part III, In Grasshopper:] Debugging script components, the canvas while paused, Panels as code
    and Markdown, the components \vn{Edit}, \vn{Text} and \vn{Protect}, and the prompts window.
  \item[Part IV, In Rhino:] Markdown as Rhino text, and links between Rhino, Grasshopper and scripts.
  \item[Part V, Reference:] The lessons, commands and keys, every setting, limits and troubleshooting.
\end{description}

The Markdown manual that \vn{Help} opens in the Script Editor (\file{miniVivisection_Manual.md}) has the
same features in short form. This manual explains them at length.

\section{Conventions}

\begin{xltabular}{\linewidth}{@{}P{0.32\linewidth}L@{}}
\toprule
\ui{Edit in Rhino}            & A button, menu item, tab or label as it appears on screen. \\
\menu{V°Watch\>Interpret as}  & A path through tabs and menus. \\
\setting{Stop once per line}  & A setting, in \vn{Settings} or \menu{Tools\>Options\>miniVivisection} (\cref{ch:settings}). \\
\keys{Ctrl+Alt+B}             & Keys to press together. \\
\code{crv.PointAtStart}       & Code: Python, a command, a value. \\
\file{examples}               & A file or folder. \\
\lessonfile{04_edit_in_rhino.py} & A lesson that teaches the feature (\cref{ch:examples}). \\
\vn{Watch}                    & One of miniVivisection's names (\see{sec:vnames}). \\
\bottomrule
\end{xltabular}
